\chapter{Jakarta e le Servlet}

\section{Jakarta EE e il Web Profile}

Nelle applicazioni web basate sul \textbf{Jakarta Web Profile} si parte dal presupposto che
sia il \emph{server HTTP} -- che esiste come applicativo indipendente -- a incapsulare
l'applicazione (modello (a) del capitolo precedente). Il server di riferimento \`e
\textbf{Apache Tomcat}, basato sul modello \emph{multi-thread}.

\begin{definizione}[Jakarta EE e Web Profile]
\textbf{Jakarta Enterprise Edition} (Jakarta EE, gi\`a \emph{Java} EE) \`e un insieme di
specifiche che estendono Java con librerie, componenti e API per lo sviluppo di applicazioni
commerciali, non incluse nel JDK standard. Il \textbf{Jakarta Web Profile} \`e il
sottoinsieme orientato allo sviluppo di applicazioni web. Tomcat implementa la specifica del
Web Profile: il suo runtime realizza un ambiente di esecuzione per l'applicazione
server-side, fatto di classi, interfacce e oggetti Java utilizzabili per realizzare le
funzionalit\`a previste.
\end{definizione}

I due strumenti principali del Jakarta Web Profile:
\begin{itemize}
  \item le \textbf{Servlet}: oggetti i cui metodi ricevono richieste HTTP, elaborano la
        risorsa richiesta e la inviano in risposta. Lo sviluppatore scrive le proprie
        classi-Servlet, istanziate da Tomcat all'avvio dell'applicazione; ogni Servlet
        specifica le route a cui risponde e Tomcat le registra nella propria tabella di
        routing, inoltrando le richieste alle servlet opportune;
  \item le \textbf{JSP} (Jakarta Server Pages): pagine in un linguaggio che intreccia HTML e
        frammenti di Java; a ciascuna Jakarta abbina una Servlet che, alla richiesta, usa la
        JSP per \emph{generare una pagina HTML} da inviare al client -- il modello delle app
        web multi-pagina \emph{orientate al server}. Nel corso non si usano: l'architettura
        scelta \`e la single-page con due applicazioni indipendenti (client-side e
        server-side) che scambiano JSON, dove la UI \`e costruita dal browser, non generata
        dal server.
\end{itemize}

\section{Architettura di esecuzione di un'app basata su Servlet}

Attori principali: il \textbf{Server HTTP} (Tomcat), l'\textbf{applicazione} (anche pi\`u
istanze sulla stessa macchina) e il \textbf{Web Container}, che fa da interfaccia fra server
e applicazioni, implementando la specifica del Web Profile e gestendo il \emph{ciclo di vita}
delle servlet.

\textbf{Il Server HTTP} riceve le richieste; rispetto alle app installate, valuta tramite il
pathname se una richiesta:
\begin{itemize}
  \item deve andare verso una delle app (tipicamente ogni app \`e caratterizzata da un
        segmento iniziale di pathname) $\rightarrow$ la passa al Web Container;
  \item va gestita altrimenti (es. servire una risorsa statica);
  \item riferisce una risorsa inesistente $\rightarrow$ risponde \texttt{404 Not Found}.
\end{itemize}

\textbf{Il Web Container} \`e responsabile di:
\begin{itemize}
  \item \emph{inizializzare} gli oggetti servlet all'avvio del server (o al primo rilascio
        dell'app);
  \item \emph{gestire le richieste} per l'applicazione e invocare i servizi della Servlet
        opportuna (anche se il pathname punta a un'app, la specifica route pu\`o non essere
        gestita: pu\`o generarsi comunque un 404);
  \item \emph{gestire le eccezioni} sorte nell'app, inviando in tal caso
        \texttt{500 Internal Server Error};
  \item \emph{distruggere} le servlet allo spegnimento del server o alla disinstallazione.
\end{itemize}

\section{Cos'\`e una Servlet: la gerarchia di classi}

\begin{lstlisting}[style=java]
package tweb.myfirstapp;
import java.io.*;
import jakarta.servlet.http.*;
import jakarta.servlet.annotation.*;

@WebServlet(name = "helloServlet", value = "/hello-servlet")
public class HelloServlet extends HttpServlet {
  private String message;
  public void init() { /* ... */ }
  public void doGet(HttpServletRequest request,
                    HttpServletResponse response) throws IOException {
    /* ... */
  }
  public void destroy() { /* ... */ }
}
\end{lstlisting}

Una Servlet \`e un'istanza di una classe che estende \texttt{HttpServlet}. La gerarchia:
\begin{itemize}
  \item l'\textbf{interfaccia \texttt{Servlet}} definisce i metodi che ogni Servlet deve
        possedere -- quelli che servono al Web Container per gestirla:
        \begin{itemize}
          \item \texttt{init(ServletConfig config)}: invocato dal container per
                inizializzare la Servlet dopo averla creata;
          \item \texttt{service(ServletRequest req, ServletResponse res)}: invocato quando il
                container chiede alla Servlet di gestire una richiesta e ottenere la
                risposta;
          \item \texttt{destroy()}: invocato prima della chiusura; la Servlet pu\`o ``fare
                pulizia'' (connessioni aperte, servizi sottoscritti, registrazioni).
        \end{itemize}
  \item la classe astratta \textbf{\texttt{GenericServlet}} implementa \texttt{Servlet} dando
        un'implementazione di base di \texttt{init(config)}: memorizza internamente
        l'oggetto config e chiama un metodo \texttt{init()} \emph{senza parametri}, che non
        fa nulla ma pu\`o essere ridefinito dalle sottoclassi;
  \item la classe \textbf{\texttt{HttpServlet}} estende GenericServlet per il caso HTTP
        (in teoria esistono Servlet per altri protocolli). Implementa il \texttt{service} di
        base, che: verifica che la richiesta sia HTTP; crea gli oggetti
        \texttt{HttpServletRequest/Response} dagli originali
        \texttt{ServletRequest/Response}; invoca una propria \texttt{service} che accetta gli
        oggetti Http; questa esamina il \emph{method HTTP} e chiama il metodo interno
        \textbf{doXXX} corrispondente (\texttt{doGet}, \texttt{doPost}, \dots). I doXXX di
        base non fanno nulla: \`e compito delle sottoclassi (le Servlet dell'applicazione)
        implementare quelli relativi alle richieste che vogliono gestire.
\end{itemize}

\section{Il ciclo di vita di una Servlet}

La vita di una Servlet inizia all'installazione (o, se gi\`a installata, all'avvio del
server) e si conclude allo shutdown del server o alla disinstallazione; nel mezzo \`e sempre
attiva e pronta a ricevere richieste. Tre momenti: \textbf{START}, \textbf{SERVICE},
\textbf{STOP}.

\begin{importante}[START]
\emph{1 -- Creazione:} il Web Container crea un \textbf{ServletContext} (uno, comune a tutta
l'applicazione); crea una \emph{singola istanza} per ciascuna Servlet dell'app; per ogni
Servlet crea un oggetto \textbf{ServletConfig} con i parametri di configurazione (usando fra
l'altro le informazioni annotate con \texttt{@WebServlet}) e un riferimento al
ServletContext; inserisce in un suo \emph{registro} le Servlet e le route a cui rispondono.
\emph{2 -- Inizializzazione:} per ogni Servlet \texttt{sN} con config \texttt{cfN}, il
container chiama \texttt{sN.init(cfN)}, che a sua volta invoca \texttt{sN.init()}. Solo al
completamento della init la Servlet \`e considerata attiva e pronta.
\end{importante}

\begin{importante}[SERVICE]
\emph{1 -- Matching \& Routing:} il container effettua un \textbf{matching} fra la URL
ricevuta e gli urlPattern dichiarati dalle Servlet, determinando la servlet ``vincente'';
l'algoritmo ha una strategia per i conflitti (es. se una URL corrisponde a due pattern,
``vince'' quello \emph{pi\`u lungo}). Se non c'\`e matching: \texttt{404 Not Found}.
\emph{2 -- Predisposizione:} vengono predisposti gli oggetti
\texttt{HttpServletRequest/Response} che estendono i \texttt{ServletRequest/Response} di
partenza (la request contiene tutte le informazioni sulla richiesta: URL, method,
header\dots).
\emph{3 -- Invocazione del servizio:} il container chiama \texttt{s.service(req, res)} di
HttpServlet con gli oggetti predisposti; questo invoca a sua volta il \texttt{doXXX}
appropriato per la richiesta in esame.
\end{importante}

\begin{importante}[STOP]
\emph{1 -- Cleaning up:} per ogni servlet il container chiama \texttt{destroy()} (di base non
fa nulla; da ridefinire se la init ha fatto cose che vanno ``disfatte'', es. chiudere una
connessione a un DB).
\emph{2 -- Garbage collection:} vengono eliminati gli oggetti Servlet, i ServletConfig e il
ServletContext.
\end{importante}
